\subsection{Native audit and performance evidence}

The driver \code{weighted\_research/suffix\_folds.py} compares against
\code{e052c7ef07f3}, the compact-coordinate implementation from the previous
section. It records 444 current source/fixture hashes and checks them before
and after each run. The baseline package is loaded from Git in a separate
namespace; only its weighted producer differs among pre-existing runtime
modules. All independent checker files are byte-for-byte unchanged.

All 1,128 maintained tests pass in 214.506 seconds. The weighted audit passes
10,000 finite systems, split evenly between the classical and Fine--Wilf
periodic rules, plus exact huge binary examples. The independent checker
passes 2,500 literal cases, 366 mutation checks, producer-disabled replay,
and a 20,002-bit endpoint case. The normal audit passes 5,100 comparisons
and certificate replays on the 1,275-vector, 45-triangulation Regina corpus.
Additional baseline comparisons give identical complete answers, including
certificates and statistics, for 2,550 normal queries and 1,000 abstract
systems. The successful audit takes 48.158 seconds while overlapping the
suite, so its duration is not a performance measurement. The outer record's
baseline and source hashes identify this run; a nested kernel record retains
its historical baseline metadata.

The first audit-driver attempt used dictionary pairings for the isolated old
producer, whose public API requires its own \code{IntervalPairing} class.
That harness error was corrected before the successful pinned audit; the
initial log is retained separately. It was not a mathematical or runtime
regression. The read-only check of \code{openai/math} found its head unchanged
at \code{fd4aeeb2ee4f}; no new upstream mathematics was imported in this step.

The complete-query benchmark runs serially after the audit, suite and first
publication. It takes 108.073 seconds and retains all 160 measured calls and
32 warm-ups. Eight supplied normal surfaces/modes use five shuffled rounds
with an A/A copy of both engines. Every timed call includes fresh geometry,
weight construction, orbit discovery, certificate production and independent
verification. Input-fixture construction, serialization and equality checks
are outside the timer. Full answers, including every certificate field and
statistic, agree on every call.

\begin{center}\small
\begin{tabular}{lrrrrr}
\toprule
Query & Old ms & New ms & Old/new & Old A/A & New A/A \\
\midrule
meridian-1 & 0.976 & 0.939 & 1.021 & 0.990 & 0.981 \\
meridian-16 & 21.757 & 21.124 & 1.038 & 1.017 & 1.012 \\
meridian-64 & 449.308 & 409.737 & 1.094 & 0.988 & 0.986 \\
meridian-128 & 2773.559 & 2462.575 & 1.126 & 1.000 & 1.001 \\
triangle-only & 4.151 & 4.172 & 1.000 & 0.999 & 1.007 \\
binary-1024 & 17.024 & 16.230 & 1.050 & 1.000 & 0.996 \\
disk-16 & 17.191 & 17.405 & 1.019 & 0.997 & 1.010 \\
disk-128 & 1213.349 & 1206.785 & 1.008 & 1.000 & 1.000 \\
\bottomrule
\end{tabular}
\end{center}


On the 128-tetrahedron full-coordinate query, median total time changes from
2.774 to 2.463 seconds, with a within-round median ratio of $1.126$.
Production medians are 1.497 and 1.225 seconds. Verification uses identical
code and proof; variations in its measured time are not an algorithmic gain.
At 64 tetrahedra the total paired ratio is $1.094$. Triangle-only queries
are unchanged ($1.000$ after rounding), while the large binary-coordinate
fixture gives $1.050$. Three-weight disk queries give only $1.019$ and
$1.008$; these small differences should not be advertised as reliable gains.
Old A/A ratios range from $0.988$ to $1.017$ and new A/A ratios from
$0.981$ to $1.012$.

This is an incremental comparison against the compact implementation, not
against the earlier dense-coordinate version. Its gains should not be
multiplied with ratios from another run to claim a measured combined speedup.
The improvement is local to weight production, leaves automatic recognition
coverage unchanged, and does not establish a general quasipolynomial bound.
All samples and deduplicated complete certificates are retained in
\code{data/suffix-fold-benchmark.json}; the table generator is
\code{data/suffix\_fold\_tables.py}.
